\section{Related Work}
\label{sec:related}

\textbf{Learned low-light raw enhancement.}
SID~\cite{chen2018sid} introduced paired short- and long-exposure raw data and a U-Net that
maps raw input straight to sRGB. ELD~\cite{wei2020eld} and LRD~\cite{zhang2023lrd} model the
physics of sensor noise to synthesise realistic training data and narrow the gap between
synthetic and real noise. These methods are regression models trained with pixel losses, so they
average over plausible outputs and oversmooth at extremely low SNR.

\textbf{Diffusion models for restoration.}
Denoising diffusion probabilistic models~\cite{ho2020ddpm} learn to reverse a gradual noising
process. DDIM~\cite{song2021ddim} samples in fewer, deterministic steps, and classifier-free
guidance (CFG)~\cite{ho2022cfg} trades fidelity against detail at inference. Latent diffusion
models~\cite{rombach2022ldm} run the process in the compressed latent space of a variational
autoencoder (VAE), which is what Stable Diffusion does. For low-light raw images,
ExposureDiffusion~\cite{wang2023exposurediffusion} trains a diffusion model from scratch on
small paired sets, which limits its generative capacity. LDM-ISP~\cite{wen2023ldmisp} adapts a
pretrained latent diffusion model with wavelet preprocessing and taming modules.
DarkDiff~\cite{zheng2025darkdiff} retasks pretrained Stable Diffusion as the ISP with three
additions: a residual VAE, region-based cross-attention to the noisy input, and a decoder image
loss. It is the starting point of this work.

\textbf{Efficient diffusion U-Nets.}
BK-SDM~\cite{kim2024bksdm} removes whole residual and attention blocks from the Stable Diffusion
U-Net and recovers quality by feature-level knowledge distillation from the full model.
Segmind's Small-SD and Tiny-SD~\cite{segmind2023tinysd} apply the same idea to a fine-tuned
SD~1.5. SDXS~\cite{song2024sdxs} distils a BK-SDM-style U-Net further, for one-step generation.
SnapFusion~\cite{li2023snapfusion} and MobileDiffusion~\cite{zhao2024mobilediffusion} design
efficient U-Nets for phones, but their weights are not public. Since DarkDiff relies on a
pretrained prior, we restrict ourselves to compact U-Nets with public weights.

\textbf{Illumination guidance.}
Retinex theory treats an image as reflectance times illumination, and many enhancement methods
estimate the illumination first. LIME~\cite{guo2017lime} refines an illumination map to brighten
an image, and RetinexNet~\cite{wei2018retinexnet} learns the decomposition.
Zero-DCE~\cite{guo2020zerodce} predicts pixel-wise brightening curves without paired data.
Retinexformer~\cite{cai2023retinexformer} uses illumination features to guide self-attention in
a transformer. We follow this line, but inject the illumination map into the cross-attention of
a pretrained diffusion U-Net, using feature-wise linear modulation (FiLM)~\cite{perez2018film}
and an attention bias.
